\section{Priorización de estudios}\label{sec:priorizacion-estudios}

Tras la selección inicial de los artículos, se procedió al análisis de los campos \textit{title}, \textit{abstract} y \textit{keywords} de cada registro. Como resultado de esta revisión, se generaron métricas de calidad para cada artículo con el objetivo de priorizar aquellos más relevantes para la investigación. Las métricas empleadas fueron las siguientes:

\begin{itemize}
	\item \textbf{SCI} \textit{(Science Citation Index)}
	\item \textbf{CVI} \textit{(Core Value Index)}
	\item \textbf{IRRQ} \textit{(Index Relation Research Question)}
\end{itemize}

Este proceso inició con un total de 156 artículos, los cuales fueron evaluados en función de su alineamiento con los objetivos de la investigación. La fase de \textit{screening} permitió descartar 139 estudios, ya que algunos abordaban herramientas tecnológicas en contextos distintos al de interés y otros no se encontraban alineados con los objetivos de investigación establecidos.
